% BEGIN AUTO-GENERATED CONTINUAL APPENDIX
\appendix
\section{Continual ASAM Pilot Study}
\label{app:continual_asam}

To assess whether the adaptive sparse-routing ideas behind ASAM can be extended to continual learning, we implemented a pilot continual-learning variant in the repository. This appendix intentionally presents the continual results as a conservative addendum rather than a replacement for the main long-sequence claims of the paper. The goal is to document the current executable continual-learning evidence, including both positive diagnostic signals and negative comparative findings.

\subsection{Setup}

We evaluate the continual extension on class-incremental Split AG News with 2 tasks of 2 classes each and a compact text classifier built from continual ASAM layers. The current paper-oriented run uses sequence length $128$, batch size $8$, $256$ training examples and $128$ validation examples per split, 1 layer with model width $64$, 4 heads, top-$k=2$ sparse pattern selection, AdamW with learning rate $3 \times 10^{-4}$, and 1 training epoch per task. Strategy-level and operator-level ablations are aggregated over 3 seeds. This setting is deliberately modest and CPU-runnable, so it should be interpreted as a controlled pilot benchmark rather than a final large-scale continual-learning claim.

\subsection{Diagnostic Benchmark}

The single-run prototype-routing benchmark with the meta-secant controller reaches average accuracy 0.4766, average forgetting 0.0312, and backward transfer -0.0312. Because this benchmark is single-seed, we use it primarily as a diagnostic case study: it is useful for inspecting stage-wise transport loss, routing stability, entropy, and prototype occupancy, but it should not be treated as the main comparative result.

\subsection{Strategy-Level Ablation}

Table~\ref{tab:continual_strategy_ablation} summarizes the current multi-seed routing/controller comparison. The strongest strategy in the present setup is no\_adaptation, which indicates that the prototype-routed continual variant can already produce measurable differences under the current pilot benchmark.

\begin{table}[htbp]
\centering
\caption{Strategy-level continual ablation on Split AG News (3 seeds). Higher accuracy and backward transfer are better; lower forgetting is better.}
\label{tab:continual_strategy_ablation}
\begin{tabular}{@{}lcccc@{}}
\toprule
\textbf{Strategy} & \textbf{Avg Acc.} & \textbf{Avg Forgetting} & \textbf{BWT} & \textbf{Final Gap} \\
\midrule
\texttt{task\_routing} & $0.5234 \pm 0.0447$ & $-0.0469 \pm 0.0776$ & $0.0469 \pm 0.0776$ & $0.0000 \pm 0.0000$ \\
\texttt{no\_adaptation} & $0.5312 \pm 0.0418$ & $-0.0625 \pm 0.0920$ & $0.0625 \pm 0.0920$ & $0.1034 \pm 0.0040$ \\
\texttt{correlation} & $0.5312 \pm 0.0418$ & $-0.0625 \pm 0.0920$ & $0.0625 \pm 0.0920$ & $0.1035 \pm 0.0028$ \\
\texttt{meta\_secant} & $0.5312 \pm 0.0418$ & $-0.0625 \pm 0.0920$ & $0.0625 \pm 0.0920$ & $0.1035 \pm 0.0028$ \\
\bottomrule
\end{tabular}
\end{table}

\subsection{Operator-Level Ablation}

Table~\ref{tab:continual_operator_ablation} isolates several operator choices inside prototype-routed continual ASAM. Task-level differences are currently modest, but the transport-facing diagnostics already reveal a meaningful pattern: \texttt{sinkhorn\_topk} and \texttt{kl\_topk} achieve the same mean accuracy, while the Sinkhorn version preserves a final transport gap of 0.1051 and the KL-based variant shows 0.0282. Disabling the transport term slightly lowers accuracy to 0.5312 and raises the final transport-loss trace to 0.0558. We therefore interpret the present operator study as an early mechanistic signal rather than a decisive empirical win.

\begin{table}[htbp]
\centering
\caption{Operator-level continual ablation on Split AG News (3 seeds).}
\label{tab:continual_operator_ablation}
\begin{tabular}{@{}lccccc@{}}
\toprule
\textbf{Operator Setting} & \textbf{Routing} & \textbf{Avg Acc.} & \textbf{Avg Forgetting} & \textbf{Final Gap} & \textbf{Final Transport} \\
\midrule
\texttt{sinkhorn\_topk} & Sinkhorn Top-$k$ & $0.5130 \pm 0.0461$ & $-0.0573 \pm 0.0940$ & $0.1051 \pm 0.0063$ & $0.0382 \pm 0.0065$ \\
\texttt{kl\_topk} & KL Top-$k$ & $0.5234 \pm 0.0522$ & $-0.0625 \pm 0.0675$ & $0.0282 \pm 0.0064$ & $0.0376 \pm 0.0060$ \\
\texttt{no\_transport} & Sinkhorn Top-$k$ & $0.5312 \pm 0.0418$ & $-0.0625 \pm 0.0920$ & $0.1034 \pm 0.0040$ & $0.0558 \pm 0.0126$ \\
\texttt{no\_merge} & Sinkhorn Top-$k$ & $0.5104 \pm 0.0266$ & $-0.0417 \pm 0.0266$ & $0.1604 \pm 0.0051$ & $0.0371 \pm 0.0059$ \\
\texttt{no\_relocation} & Sinkhorn Top-$k$ & $0.5130 \pm 0.0461$ & $-0.0573 \pm 0.0940$ & $0.1166 \pm 0.0221$ & $0.0373 \pm 0.0064$ \\
\bottomrule
\end{tabular}
\end{table}

\subsection{Takeaway}

The current continual-learning results support a cautious positive pilot conclusion. The continual ASAM framework already provides executable transport-aware diagnostics, prototype lifecycle statistics, and operator-level ablations. In this specific small-scale CPU run, the best prototype-routed variant improves average accuracy over \texttt{task\_routing} by 0.0078 and reduces average forgetting by 0.0156. At the same time, the margin is small and the benchmark remains low-budget (3 seeds with modest sample counts), so this should be treated as an encouraging pilot signal rather than a decisive retention claim. Stronger data scale, longer training, and richer ablations are still needed before claiming a robust empirical advantage.
% END AUTO-GENERATED CONTINUAL APPENDIX
